% GENERATED FILE. Edit canonical lessons, then run npm run generate:books.
% canonical-source-hash: fnv1a64:4f8d9f93c516ff2a
% canonical-lessons: MR-C444-dabke, MR-C444-gal, MR-C444-gadhul, MR-C444-jhirapne, MR-C444-osandne

\chapter{Puddles and Seepage}
\label{ch:mr-puddles-and-seepage}
\hlchaptermodality{444}

\begin{chapteropening}
\textbf{By the end of this chapter:} \emph{I can talk about puddles, silt and muddy water, and say water seeped or overflowed.}

Complete the last lesson of chapter 444: I can talk about puddles, silt and muddy water, and say water seeped or overflowed.
\end{chapteropening}

\section[\texorpdfstring{\mr{डबके} (ḍabke)}{डबके (ḍabke)}]{\mr{डबके} (ḍabke) — a puddle}

\label{lesson:MR-C444-dabke}



\begin{quote}
Before the new one: say the Marathi for a large lake, then the Marathi for to gurgle.
\end{quote}

\subsection*{You'll want to know: \mr{डबके}}
\textbf{\mr{डबके}} — \emph{ḍabke} — \textquotedblleft{}a puddle\textquotedblright{}.

\begin{quote}
\mr{रस्त्यावर} \mr{पावसाचं} \mr{डबकं} \mr{साचलं}.
\end{quote}

\begin{quote}
\emph{rastyāvar pāvsācaṁ ḍabkaṁ sāclaṁ.}
\end{quote}

\begin{quote}
A puddle of rain collected on the road.
\end{quote}

Neuter. Plural: \mr{डबकी}, \emph{ḍabkī}.

Written \mr{डबके}, said and often written \mr{डबकं}; the plural is \mr{डबकी}.

\begin{cousinweb}
\textbf{\mr{डबके}} (\emph{ḍabke}) is a Marathi word for a small pool of standing water; its origin is not settled.
\end{cousinweb}

\begin{grammarlens}[title={What you've built}]
One more word for getting around.
\end{grammarlens}

\subsection*{Your turn}
\begin{itemize}
\raggedright
  \item \emph{Say it:} \emph{ḍabke}
  \item \emph{Say it:} \emph{ḍabke}, once more
  \item \emph{Say it:} \emph{rastyāvar pāvsācaṁ ḍabkaṁ sāclaṁ.} — A puddle of rain collected on the road.
  \item \emph{Recall:} say the Marathi for a large lake, then the Marathi for to gurgle, then say \emph{ḍabke} again
\end{itemize}

\begin{tcolorbox}[breakable,colback=teal!4,colframe=teal!35!black,title={Before you move on}]
What does \mr{डबके} mean? (\textquotedblleft{}A puddle\textquotedblright{}.) Say it once more.
\end{tcolorbox}
\section[\texorpdfstring{\mr{गाळ} (gāḷ)}{गाळ (gāḷ)}]{\mr{गाळ} (gāḷ) — silt}

\label{lesson:MR-C444-gal}



\begin{quote}
Before the new one: say the Marathi for to gurgle, then the Marathi for a puddle.
\end{quote}

\subsection*{You'll want to know: \mr{गाळ}}
\textbf{\mr{गाळ}} — \emph{gāḷ} — \textquotedblleft{}silt\textquotedblright{}.

\begin{quote}
\mr{पुरानंतर} \mr{शेतात} \mr{गाळ} \mr{साचला}.
\end{quote}

\begin{quote}
\emph{purānantar śetāt gāḷ sāclā.}
\end{quote}

\begin{quote}
After the flood silt settled in the fields.
\end{quote}

Masculine.

\mr{गाळ} \mr{काढणे}, to clear the silt from a tank or a well.

\begin{cousinweb}
\textbf{\mr{गाळ}} (\emph{gāḷ}) is inherited: CDIAL lists Marathi \emph{gāḷ}, \textquotedblleft{}dregs, mud\textquotedblright{}, under \emph{*gāḍa} (CDIAL 4114).
\end{cousinweb}

\begin{grammarlens}[title={What you've built}]
One more word for getting around.
\end{grammarlens}

\subsection*{Your turn}
\begin{itemize}
\raggedright
  \item \emph{Say it:} \emph{gāḷ}
  \item \emph{Say it:} \emph{gāḷ}, once more
  \item \emph{Say it:} \emph{purānantar śetāt gāḷ sāclā.} — After the flood silt settled in the fields.
  \item \emph{Recall:} say the Marathi for to gurgle, then the Marathi for a puddle, then say \emph{gāḷ} again
\end{itemize}

\begin{tcolorbox}[breakable,colback=teal!4,colframe=teal!35!black,title={Before you move on}]
What does \mr{गाळ} mean? (\textquotedblleft{}Silt\textquotedblright{}.) Say it once more.
\end{tcolorbox}
\section[\texorpdfstring{\mr{गढूळ} (gaḍhūḷ)}{गढूळ (gaḍhūḷ)}]{\mr{गढूळ} (gaḍhūḷ) — muddy; murky}

\label{lesson:MR-C444-gadhul}



\begin{quote}
Before the new one: say the Marathi for a puddle, then the Marathi for silt.
\end{quote}

\subsection*{You'll want to know: \mr{गढूळ}}
\textbf{\mr{गढूळ}} — \emph{gaḍhūḷ} — \textquotedblleft{}muddy; murky\textquotedblright{}.

\begin{quote}
\mr{पावसाळ्यात} \mr{नदीचं} \mr{पाणी} \mr{गढूळ} \mr{होतं}.
\end{quote}

\begin{quote}
\emph{pāvsāḷyāt nadīcaṁ pāṇī gaḍhūḷ hotaṁ.}
\end{quote}

\begin{quote}
In the rainy season the river water turns muddy.
\end{quote}

It does not change: \mr{गढूळ} \mr{पाणी}, \mr{गढूळ} \mr{नदी}.

\begin{cousinweb}
\textbf{\mr{गढूळ}} (\emph{gaḍhūḷ}) is a Marathi adjective for cloudy, muddy water; its origin is not settled.
\end{cousinweb}

\begin{grammarlens}[title={What you've built}]
One more word for getting around.
\end{grammarlens}

\subsection*{Your turn}
\begin{itemize}
\raggedright
  \item \emph{Say it:} \emph{gaḍhūḷ}
  \item \emph{Say it:} \emph{gaḍhūḷ}, once more
  \item \emph{Say it:} \emph{pāvsāḷyāt nadīcaṁ pāṇī gaḍhūḷ hotaṁ.} — In the rainy season the river water turns muddy.
  \item \emph{Recall:} say the Marathi for a puddle, then the Marathi for silt, then say \emph{gaḍhūḷ} again
\end{itemize}

\begin{tcolorbox}[breakable,colback=teal!4,colframe=teal!35!black,title={Before you move on}]
What does \mr{गढूळ} mean? (\textquotedblleft{}Muddy; murky\textquotedblright{}.) Say it once more.
\end{tcolorbox}
\section[\texorpdfstring{\mr{झिरपणे} (jhirapṇe)}{झिरपणे (jhirapṇe)}]{\mr{झिरपणे} (jhirapṇe) — to seep}

\label{lesson:MR-C444-jhirapne}



\begin{quote}
Before the new one: say the Marathi for silt, then the Marathi for muddy.
\end{quote}

\subsection*{You'll want to know: \mr{झिरपणे}}
\textbf{\mr{झिरपणे}} — \emph{jhirapṇe} — \textquotedblleft{}to seep\textquotedblright{}.

\begin{quote}
\mr{पावसाचं} \mr{पाणी} \mr{छपरातून} \mr{झिरपलं}.
\end{quote}

\begin{quote}
\emph{pāvsācaṁ pāṇī chaprātūn jhiraplaṁ.}
\end{quote}

\begin{quote}
The rainwater seeped through the roof.
\end{quote}

Past (he …): \mr{तो} \mr{झिरपला}, \emph{to jhiraplā}.

Slower than \mr{पाझरणे}: drop by drop.

\begin{cousinweb}
\textbf{\mr{झिरपणे}} (\emph{jhirapṇe}) is inherited: CDIAL lists Marathi \emph{jhirapṇẽ}, \textquotedblleft{}to trickle\textquotedblright{}, under \emph{*jhirati} (CDIAL 5390).
\end{cousinweb}

\begin{grammarlens}[title={What you've built}]
One more word for getting around.
\end{grammarlens}

\subsection*{Your turn}
\begin{itemize}
\raggedright
  \item \emph{Say it:} \emph{jhirapṇe}
  \item \emph{Say it:} \emph{jhirapṇe}, once more
  \item \emph{Say it:} \emph{pāvsācaṁ pāṇī chaprātūn jhiraplaṁ.} — The rainwater seeped through the roof.
  \item \emph{Recall:} say the Marathi for silt, then the Marathi for muddy, then say \emph{jhirapṇe} again
\end{itemize}

\begin{tcolorbox}[breakable,colback=teal!4,colframe=teal!35!black,title={Before you move on}]
What does \mr{झिरपणे} mean? (\textquotedblleft{}To seep\textquotedblright{}.) Say it once more.
\end{tcolorbox}
\section[\texorpdfstring{\mr{ओसंडणे} (osaṇḍṇe)}{ओसंडणे (osaṇḍṇe)}]{\mr{ओसंडणे} (osaṇḍṇe) — to overflow}

\label{lesson:MR-C444-osandne}



\begin{quote}
Before the new one: say the Marathi for muddy, then the Marathi for to seep.
\end{quote}

\subsection*{You'll want to know: \mr{ओसंडणे}}
\textbf{\mr{ओसंडणे}} — \emph{osaṇḍṇe} — \textquotedblleft{}to overflow\textquotedblright{}.

\begin{quote}
\mr{धरण} \mr{भरून} \mr{ओसंडलं}.
\end{quote}

\begin{quote}
\emph{dharaṇ bharūn osaṇḍlaṁ.}
\end{quote}

\begin{quote}
The dam filled and overflowed.
\end{quote}

Past (he …): \mr{तो} \mr{ओसंडला}, \emph{to osaṇḍlā}.

Often with \mr{भरून}: \mr{भरून} \mr{ओसंडणे}, to be full to overflowing.

\begin{cousinweb}
\textbf{\mr{ओसंडणे}} (\emph{osaṇḍṇe}) is a Marathi verb for water spilling over the top; its origin is not settled.
\end{cousinweb}

\begin{grammarlens}[title={What you've built}]
One more word for getting around.
\end{grammarlens}

\subsection*{Your turn}
\begin{itemize}
\raggedright
  \item \emph{Say it:} \emph{osaṇḍṇe}
  \item \emph{Say it:} \emph{osaṇḍṇe}, once more
  \item \emph{Say it:} \emph{dharaṇ bharūn osaṇḍlaṁ.} — The dam filled and overflowed.
  \item \emph{Recall:} say the Marathi for muddy, then the Marathi for to seep, then say \emph{osaṇḍṇe} again
\end{itemize}

\begin{tcolorbox}[breakable,colback=teal!4,colframe=teal!35!black,title={Before you move on}]
What does \mr{ओसंडणे} mean? (\textquotedblleft{}To overflow\textquotedblright{}.) Say it once more.
\end{tcolorbox}
